\section{项目背景与需求分析}

\subsection{项目背景}
选题来自学生在同一时段内处理多类校园事务的经历。以准备次日上课为例，学生先查教务课表确认节次和教室，再到学习通检查作业、考试或活动，最后把平台之外的个人事项另记一份。每项操作都已有服务承接，真正重复的是寻找入口、判断消息归属和确认上次查询是否仍然有效。校园饮水和洗浴又各自依赖不同账户。分散带来的困难有两层。学生先要判断该从哪个入口查起，随后还要判断看到的是平台刚返回的结果、上次保存的内容，还是一次查询失败后的空白。以跨周调课为例，知道课程名称并不足以确定走课路线，还需要核对周次与地点。以提交作业为例，看到课堂活动的标题也不能直接判断任务是否仍待处理。酱菜通由此以校园助手为主题，将分散的服务入口集中起来，并清楚呈现查询结果及其状态。

\subsection{开发动机与目标}
本项目将学生每天重复打开的入口收进一个应用。课表适合作为启动后的首屏，进入首页即读本机最近成功数据，随后按条件同步。学习页以任务和教学班为组织单位，提供类型、状态及课程分组，便于在课前判断哪些工作仍待处理。待办页负责个人自行添加的事项，网络不可用时仍可编辑和恢复。生活页集中饮水与洗浴入口，并以服务实际接入状态决定可用操作。我将整合范围限定在共同的入口和操作习惯。课表沿着学期、周次和星期查看，学习页按教学班和任务状态收拢信息，个人待办独立编辑，生活数据继续以各服务的账号和回执为准。

\subsection{目标用户与主要需求}
本项目主要面向江西财经大学在校学生，围绕课程查询、学习任务、个人待办及校园生活服务等日常需求展开设计。走课场景关注课程时间与地点，要求课表在网络波动时仍可查阅。学习管理场景关注是否待交、何时截止，以及哪个教学班发布了任务，过期记录仍应可追溯。个人待办要求保存简单、完成状态可反转。余额和账单等生活数据则需要从本人会话读取，账号失效后应有明确重连路径。这些场景分别对应三项设计判断。课表是走课时需要立即查到的信息，因此启动优先回读，保留更新状态；学习通和待办都可以用状态筛选，但前者采用平台任务状态，后者的完成标记由本人编辑。余额及交易依赖外部会话，因此用户可以集中查看连接状态，并在准备好时主动重连。

\subsection{功能范围}
项目聚焦校园信息查询、任务组织和本机操作。学校成绩和考试安排以校方查询结果为准，学习通开放任务跳转，独立待办由本机保存。慧生活798提供独立的设备和积分流程，实际出水或积分扣减等动作须由用户明确操作。住理生活的原生接入仍处于部分实现，服务页显示设备接口状态并提供官方入口。项目以现有校方和第三方服务提供动态数据，应用直接在本机完成整理和缓存。
